
\section{Hyperspherical Parent and Harmonic Convention}

We use a closed FLRW geometry
\[
ds^2=-dt^2+a^2(t)R_H^2\left[d\psi^2+\sin^2\psi\,d\Omega_2^2\right],
\]
with \(a_0=1\). Here \(\psi\) is a dimensionless radial angle on the unit three-sphere, \(R_H\) is the present-day comoving curvature radius, and the physical curvature radius is \(aR_H\). To avoid mixing dimensional and angular radial coordinates, define
\[
D_C(z)=c\int_0^z\frac{dz'}{H(z')},
\qquad
\psi(z)=\frac{D_C(z)}{R_H}.
\]

Scalar harmonics satisfy
\[
-\nabla_\gamma^2Y_n=\frac{n(n+2)}{R_H^2}Y_n,
\qquad
n=0,1,2,\ldots,
\]
with degeneracy \((n+1)^2\). The physical spatial eigenvalue entering the time-dependent perturbation equations is
\[
p_n(a)=\frac{n(n+2)}{a^2R_H^2}.
\]

The first nonconstant eigenspace is therefore
\[
p_1=\frac{3}{a^2R_H^2}.
\]
This convention is used throughout. Older scalar-topographic notes employed more than one normalization for the lowest mode; no numerical forecast is carried forward here without first converting it to the convention above.

\subsection{Exact observer-sky form of the n=1 scalar mode}

Embed the unit three-sphere in \(\mathbb R^4\). The \(n=1\) scalar harmonics are the four coordinate functions restricted to \(S^3\). A general real \(n=1\) scalar mode can therefore be written
\[
T_1(\psi,\hat{\mathbf n})
=
A\left[
\cos\alpha\cos\psi
+
\sin\alpha\sin\psi\,
(\hat{\mathbf n}\cdot\hat{\mathbf p})
\right].
\]
At fixed radial shell this scalar function contains an exact sky monopole plus an exact sky dipole and no higher angular multipoles. Its dipolar part is
\[
T_{1,{\rm dip}}(z,\hat{\mathbf n})
=
A\sin\alpha\,
\sin\!\psi(z)\,
(\hat{\mathbf n}\cdot\hat{\mathbf p}).
\]
This removes an ambiguity present in simplified radial descriptions: the zonal term \(\cos\psi\) alone is isotropic for an observer at its pole. A sky dipole in the scalar field is supplied by a generic orientation of the full \(n=1\) eigenspace relative to the observer.

\subsection{Exceptional lowest scalar harmonic and metric-response caveat}

The existence of the scalar function above must be distinguished from its use as an ordinary scalar \emph{metric} perturbation of closed FLRW. A common closed-universe convention labels scalar eigenvalues as
\[
k^2=N^2-1,
\qquad N=1,2,3,\ldots
\]
for unit positive curvature. Our convention \(n(n+2)\) is related by
\[
N=n+1.
\]
Thus our first nonconstant \(n=1\) scalar level maps to \(N=2\) in that literature.

Gauge-invariant closed-FLRW perturbation constructions contain factors such as
\[
\frac{1}{N^2-4\mathcal K},
\]
which are singular at \(N=2\) for \(\mathcal K=+1\), and the usual scalar-derived tensor-harmonic construction treats the lowest \(N=1,2\) sectors as exceptional \citep{KieferVardanyan2022}. This does not remove the \(S^3\) scalar eigenfunction itself. It means that the ordinary metric-perturbation construction degenerates precisely at the harmonic level used here.

Accordingly, the monopole-plus-dipole result above is an exact statement about the topographic scalar pattern. A physical luminosity-distance, velocity, or gravitational-potential dipole does \emph{not} follow from that geometry alone. The bridge requires a separate constrained, gauge-invariant derivation of the scalar--metric response in this exceptional sector. Until that calculation is supplied, every metric or light-cone prediction sourced by the \(n=1\) field is conditional on this consistency requirement.
